\section{n3. Continuous geometric envelopes on normal finite-group quotients in all dimensions}
\textit{Supports: ENTIRE qualified G331; n1; n2.}

\begin{theorem}
Let $k$ be algebraically closed of characteristic zero, let $B$ be a smooth
integral affine $k$-algebra, and let a finite group $G$ act by $k$-algebra
automorphisms fixing a closed point $\mathfrak p$. Set
$R=(B^G)_{\mathfrak p\cap B^G}$ and $S=B_{\mathfrak p}$.
Alternatively let $S=k[[x_1,\ldots,x_d]]$ with ANY finite group of continuous
local $k$-algebra automorphisms, and set $R=S^G$.
Then $R$ is an excellent NORMAL local domain, $S$ is finite over $R$, and
for EVERY primary $\mathfrak a\subset R$ and EVERY finite continuous signed
obstacle $f$ on the full $V_{\mathfrak a}(R)$, the independently geometric
$C_R$ and $D_R$ minorants have the same greatest member $P_R f\in C_R$.
It satisfies
\[
 (P_R f)\circ\rho=P_S(f\circ\rho),\qquad
 \rho:V_{\mathfrak aS}(S)\longrightarrow V_{\mathfrak a}(R).
\]
Every canonical degree $h_m^R$ converges uniformly to $P_R f$ with actual
negative effective ample positivity before limits. For a finite signed
primary model obstacle the target error is $O(1/m)$ on EVERY support.
\end{theorem}

Here are the invariant-ring hypotheses, proved rather than assumed.
In the affine case choose algebra generators of $B$ and let $A_0$ be the
finitely generated $k$-algebra of coefficients of their monic orbit
polynomials. Each generator is integral over $A_0$, so $B$ is a finite
$A_0$-module. The $A_0$-submodule $B^G\subset B$ is finite because $A_0$
is Noetherian. Thus $B^G$ is finitely generated over $k$ and $B$ is finite
over it. If an element of its fraction field is integral over $B^G$, it
belongs to normal $B$ and is fixed by $G$, so belongs to $B^G$. Hence
$B^G$ and $R$ are normal, and finite-type field algebras and their
localizations are excellent.

There is just one prime of $B$ above the marked maximal ideal
$\mathfrak q=\mathfrak p\cap B^G$. Any such prime $P$ is maximal. If
$P\ne\mathfrak p$, choose $b\in P\setminus\mathfrak p$. Since the point
is $G$ fixed, every $\sigma(b)$ is outside $\mathfrak p$. But their
product belongs to $P\cap B^G=\mathfrak q\subset\mathfrak p$, a
contradiction. Therefore the finite localization $B\otimes_{B^G}R$
is local with maximal ideal $\mathfrak p$, and equals $B_{\mathfrak p}$:
every element outside that ideal is a unit. Since its denominators from
the base are invariant, taking invariants in this localization gives
$S^G=R$. Dimensions agree by finite integrality. Thus the source is
exactly a smooth affine CLOSED germ at the BFJ hypotheses.

For the formal case write $d_G=|G|$. The coefficients of the orbit
polynomials of the $x_i$ lie in the maximal ideal of $S$. Map a power
series ring $k[[z_1,\ldots,z_l]]$ to $S$ by sending its variables to those
coefficients, and let $A_0$ denote the quotient by the kernel, identified
with its image. This is a complete Noetherian local ring. Its coefficient
ideal extended to $S$ contains every $x_i^{d_G}$, by the orbit equations,
so has maximal-ideal radical. Choose finitely many monomials spanning
the Artinian quotient of $S$ by that extended ideal. They generate $S$
over $A_0$: subtract their $k$-linear residue combination, write the
remainder in the coefficient ideal, and repeat. At step $N$ the residual
lies in the $N$th power of this ideal and hence in $\mathfrak m_S^N$.
The coefficients of the chosen monomials are convergent formal series in
the $z_j$; their images yield the desired expression in the complete ring
$S$. This proves finiteness without assuming algebraizability of the action.

Now $R=S^G$ is a finite $A_0$-submodule of $S$, hence Noetherian, and $S$
is finite over $R$. It is local: an invariant element outside
$\mathfrak m_S$ has invariant inverse. It is complete as a finite module
over complete $A_0$; the extended maximal ideal of $A_0$ is primary for
its unique maximal ideal, so this is also maximal-ideal completeness.
Complete Noetherian local rings are excellent, the same ordinary
excellence background retained in G331. Normality follows as in the
affine case. The source is the completion of the smooth polynomial germ;
G331's full boundary homeomorphism and exact primary contraction transfer
n1's CE and model rate to it. This does not algebraize the finite action
or assert equivalence of arbitrary geometric model categories.

Apply n2: EVERY full restriction fibre is a $G$ orbit, including all
supports. The pulled-back obstacle $F=f\circ\rho$ is invariant, and group
automorphisms permute the geometric classes (actual models, effectivity
and curve degrees are preserved; equivalently use G331's comparison).
The unique continuous greatest source envelope $H=P_S F$ from n1 is
therefore invariant. The ENTIRE G331 full-fibre POSITIVE descent theorem
gives $h\in C_R$ with $h\circ\rho=H$. It is below $f$. Every target
$C/D$ competitor pulls back to a source competitor and is below $H$,
so it is below $h$ by surjectivity. This proves the theorem, including
$D$ competitors with partial bottom. The target is NORMAL; this is a
new higher-dimensional normal sector, not a conductor counterexample.
Its normalization is itself, with unit conductor and no infinite-conductor
normalization stratum. The nontrivial finite-cover fibres are instead
exactly the $G$ orbits proved in n2.

Positivity can also be seen through explicit actual profiles. Approximate
invariant $H$ uniformly by a rational source primary/unit profile
$-w(J)/n$. Average its translates. The orbit coefficient ideal of n2
gives the downstairs profile $-v(\mathcal N_G(J))/(d_G n)$ with the same
error. Its own ACTUAL primary blowup has negative effective ample
Cartier divisor $-E/(d_G n)$, positive on EVERY contracted integral curve
before the limit. Thus no finite-fibre minimum is declared psh and no
numerical divisor on an unavailable model is used.

Here is the all-degree quantitative assertion. Let $J_m=A_m^S(F)$;
these ideals are $G$ stable because $F$ is invariant. The exact canonical
contraction is
\[
 A_m^R(f)=J_m\cap R.
\]
It follows directly from surjectivity and equality of every restricted
value; equality after EXTENDING this contraction is not asserted.
For $d_G=d$, n2 gives
\[
 \mathcal N_G(J_m)\subset A_{dm}^R(f),\qquad
 h_{dm}^R\circ\rho\ge -w(J_m)/m=h_m^S.
\]
Indeed its coefficient value is $d\,w(J_m)\ge-dm f(v)$ on EVERY base
point, so every coefficient satisfies the defining canonical tests.
For a signed model obstacle, n1 yields $0\le H-h_m^S\le K/m$ with a
finite $K$ at the transported primary gauge. Let
$M_1=\|v(A_1^R(f))\|_\infty$. For $l=dm+s$, $0\le s<d$, $m\ge1$,
canonical grading gives
\[
 h_l^R\ge(dm/l)h_{dm}^R-sM_1/l
       \ge P_Rf-(dK+sM_1)/l.
\]
Here $P_R f\le0$, so the missing $s/l$ fraction of it improves the
inequality. For $1\le l<d$, $(A_1)^l\subset A_l$ gives $h_l\ge-M_1$.
Consequently the single estimate
\[
 0\le P_Rf-h_l^R\le d(K+M_1)/l\quad(l\ge1)
\]
holds on ALL supports. If $f\ge0$ all ideals are units and the error is
zero. For an arbitrary continuous obstacle only uniform convergence is
claimed, supplied by n1 once continuity of the target envelope is proved.
Finite generation, a rational envelope or an algebraic rate for arbitrary
continuous data does not follow.

The envelope map is monotone and 1-Lipschitz in the chosen supremum norm.
A decreasing net of finite continuous obstacles with these $C$ envelopes
has a greatest $D$ minorant given by their actual decreasing envelope net:
it is below the limit, and every competitor is below every term. The
identically bottom limit is permitted. There is no arbitrary-net
flattening assumption. Exact primary changes and completion obey n1/G331;
constant denominators on finite restriction fibres give the stated
pullback formula under every primary gauge.

For an explicit genuinely singular normal example in dimension $d\ge3$,
take the involution changing signs of $x_1,x_2$ and fixing $x_3,\ldots,x_d$.
Then
\[
 R=k[a,b,c,x_3,\ldots,x_d]_0/(ac-b^2),\quad
 a=x_1^2, b=x_1x_2, c=x_2^2.
\]
Even total degree monomials in $x_1,x_2$ are generated by $a,b,c$;
the relation is $ac=b^2$, and reducing its power of $b$ to zero or one
gives distinct monomials upstairs, proving that it is the whole kernel.
The invariant ring is normal by the proof above, and the marked point is
singular since its cotangent dimension is $d+1$ while its dimension is $d$.
The theorem applies to every continuous obstacle there, and also to its
completion. This example is illustrative; the theorem treats all the
finite actions stated above.

In arbitrary positive or mixed characteristic, n2 and the conditional
transfer proof remain valid when the ring hypotheses and source CE are
actually available. We do NOT infer that source CE in higher dimension.
No arbitrary excellent singular germ is presumed to admit a finite regular
cover. General excellent normal dimension at least three, unrestricted
conductor gluing, products, energy and full P12.6 operators remain residual.

Artifacts: notes/n3.tex, P12.3-quotients-and-wild-conductors.tex and the
separate verdict/ledger. Rerun: python exact\_checks.py;
python build\_delivery.py; python verify\_delivery.py. This is an unverified
complete proof of the stated quotient scope, not original P12.3 closure.
